\section{Balanced-Force 条件：離散化演算子不一致の誤差解析}
\label{app:balanced_force_taylor}
% ═══════════════════════════════════════════════════════════════════════════════

本節では，§\ref{sec:bf_operator_mismatch}（式~\eqref{eq:bf_operator_mismatch}）で示した演算子不一致誤差を，離散テイラー展開によって定量的に導出する．

\subsection{打ち切り誤差の定量評価}

連続レベルでは $\bnabla p = \kappa\bnabla\psi/We$ の釣り合いが成立している静止液滴を，
異なる差分演算子 $D^{(1)}_{\mathrm{low}}$（低次）と $D^{(1)}_{\mathrm{CCD}}$（CCD，高次）で計算した場合の離散誤差を見る．

$f$ が解析関数のとき，各演算子の打ち切り誤差は：
\begin{align}
  D^{(1)}_{\mathrm{low}} f   &= f' + c_1 h^2 f''' + \Ord{h^4} \\
  D^{(1)}_{\mathrm{CCD}} f   &= f' + c_2 h^6 f^{(7)} + \Ord{h^8}
\end{align}

\textbf{平衡条件 $\nabla p = (\kappa/We)\nabla\psi$ を用いた残差の展開：}
\begin{align}
  D^{(1)}_{\mathrm{CCD}} p - \frac{\kappa}{We} D^{(1)}_{\mathrm{low}}\psi
  &= \Bigl[\nabla p + \Ord{h^6}\Bigr]
     - \frac{\kappa}{We}\Bigl[\nabla\psi + \Ord{h^2}\Bigr] \notag\\
  &= \underbrace{\Bigl(\nabla p - \frac{\kappa}{We}\nabla\psi\Bigr)}_{=\,0\;\text{（平衡条件）}}
     + \Ord{h^6} - \frac{\kappa}{We}\Ord{h^2} \notag\\
  &= \Ord{h^2}
  \label{eq:bf_mismatch_expand}
\end{align}
FD の打ち切り誤差 $O(h^2)$ が残差を支配し（CCD の $O(h^6)$ は無視できる），この $O(h^2)$ の不整合が\textbf{体積力として運動量方程式に残り，寄生流れを駆動し，最終的に計算の爆発を引き起こす}（§\ref{sec:bf_operator_mismatch}）．

\subsection{同一演算子（CCD）使用時の誤差低減}

同一の演算子（CCD）を両方に用いると，各項が独立に $O(h^6)$ 精度で連続微分を近似するため（$D_{\mathrm{CCD}}^{(1)} f = \nabla f + O(h^6)$），平衡条件 $\nabla p = (\kappa/We)\nabla\psi$ のもとで：
\[
  D^{(1)}_{\mathrm{CCD}} p - \frac{\kappa}{We} D^{(1)}_{\mathrm{CCD}}\psi
  = \underbrace{\nabla p - \frac{\kappa}{We}\nabla\psi}_{=\,0\;\text{（平衡条件）}} + O(h^6)
  = O(h^6)
\]
となり，寄生流れの\textbf{数値離散化誤差成分}は $O(h^6)$ まで低減される．

\noindent\textbf{（注：演算子の可換性について）}
$D_{\mathrm{CCD}}(\kappa\psi/We) = (\kappa/We)D_{\mathrm{CCD}}\psi + (\psi/We)D_{\mathrm{CCD}}\kappa$
であるから，$\kappa$ が空間変化する一般の場合は代数的因数分解 $D_{\mathrm{CCD}}(p-\kappa\psi/We)$ は使えない．上式の $O(h^6)$ は各演算子の精度から直接導かれる．

\subsection{CSF モデル誤差との関係：実測収束次数}

ただし，残差は CSF モデルが界面を有限厚さ $\varepsilon \sim O(h)$ で近似することによる\textbf{モデル誤差 $O(\varepsilon^2) \approx O(h^2)$} が律速するため，実測の寄生流れ収束次数は $O(h^2)$ 程度となる（§\ref{sec:accuracy_summary} および第\ref{sec:verification}章の静止液滴ベンチマーク参照）．Balanced-Force 条件の「Balanced」とは「数値離散化誤差成分の相殺（$O(h^6)$ への低減）」を指し，CSF モデル誤差（$O(\varepsilon^2) \approx O(h^2)$）そのものを除去するものではない点に注意されたい．

Rhie--Chow 補正と Balanced-Force 条件の独立性は付録~\ref{app:rc_precision} 末尾で確認されている（表面張力補正項を含む拡張形式も同様）．

% ═══════════════════════════════════════════════════════════════════════════════
